\section{RSM 通量与采样范围修正}\label{sec:rsm-energy}
经典场景仍出现不合理的间接光后，继续调强度难以判断问题来自哪里。前轮测试已经证明 RSM 附件接通、普通与实例化都读取它，但这些检查没有证明反射光的数值正确。本次先拆开通量写入、间接光收集和接收材质，分别与解析结果比较。

\subsection{从距离衰减排查到面积缺失}
原论文的几何因子分母是距离四次方，项目则是距离平方。逐项展开后发现，论文的两个点积使用未单位化的位移，各带一次距离。项目先单位化，因此这两种写法等价。距离平方本身不是错误。

真正不一致的是 RSMDepth 把底色、光强与入射余弦相乘，称为通量，收集阶段又只对样本求平均。方向光的一个纹素代表固定投影面积，倾斜表面的实际面积增长恰好抵消入射余弦。旧实现既多乘了该余弦，也没有把采样圆盘面积计入结果。强度倍率会掩盖其中一部分亮度差，无法恢复正确的尺度关系。

现在附件保存按光源投影面积计的反射通量密度，收集时统一乘圆盘投影面积，并计入发送、接收两端的 Lambert 系数。收发材质都使用 $(1-\text{metallic})\,\text{albedo}$，使纯金属不再参与这项漫反射。完整推导与旧提交 \code{1a168c3} 的最短代码对照放在《渲染公式与实现》的 RSM 节。

\subsection{采样数量不再缩小覆盖范围}
原始圆盘按 128 个点生成，却只取用户请求的前 N 个。前面的点更靠中心，减少到 16 个时覆盖半径只剩约 35\%。现在按实际 N 生成等面积圆盘，参数变化时清除旧历史，避免继续混入旧设置的结果。

固定的距离偏移和贡献上限改为一个光源纹素宽度的最小距离。它随场景覆盖尺度变化，但仍是近距离近似。圆盘外或没有表面的样本贡献为零，分母保留请求样本数，不按有效数量重新放大。

\subsection{实际验证与剩余限制}
Renderer3DTests 分别检查圆盘面积、距离翻倍的衰减、相同世界采样半径的结果、正照与斜照，以及背面和纯金属发送端。DepthPrepassTests 检查纯金属接收端。两组真实 GPU 数值测试通过，具体输入与误差列在公式文档。

两反射点之间仍没有可见性查询，挡板后的漏光属于尚未实现的遮挡处理。当前只有一次间接漫反射，历史也没有物体运动与前帧深度拒绝。新公式通过不能扩大成复杂场景全部正确，经典场景的图片与对照另见《经典场景实测》。

\loc{RedflagForge/assets/shaders/RSMDepth.glsl}
\loc{RedflagForge/assets/shaders/RSMIndirect.glsl}
\loc{Redflag/src/Redflag/Renderer/SceneRenderer.cpp}
\note{依据：\source{https://users.soe.ucsc.edu/~pang/160/s13/proposal/mijallen/proposal/media/p203-dachsbacher.pdf}{Dachsbacher、Stamminger：Reflective Shadow Maps，§3，印刷页 204--205}。验证记录见 \path{output/validation/rsm-regression.txt}，场景对照在 \path{output/scenes/rsm/}。旧 Mask 的 RenderDoc 捕获不用于证明此次 RSM 数值。}
